\section{Nguồn kết quả và cách tái tính}


Các bảng trong chương này được rút từ kết quả đã có; lần biên tập không chạy
thêm thí nghiệm để thay một kết quả bất lợi. Bảng~\ref{bang:nguonketqua} chỉ
rõ hồ sơ và tệp phân tích cho ba lớp bằng chứng trung tâm. Tên tệp được giữ
để có thể tìm đúng nguồn trong gói nghiên cứu, thay vì chỉ tìm một bảng đã
làm tròn trong PDF.

\begin{table}[htbp]
\centering
\caption{Hồ sơ truy nguồn cho kết quả trung tâm. Các tệp kiểm lại là phân tích
đọc dữ liệu đã ghi, không phải bộ chạy mô phỏng khác.}
\label{bang:nguonketqua}
\small
\setlength{\tabcolsep}{5pt}
\begin{tabularx}{\linewidth}{@{}p{2cm}XX@{}}
\toprule
\textbf{Nhóm} & \textbf{Hồ sơ kết quả hoặc kế hoạch} & \textbf{Tệp phân tích và đối chiếu}\\
\midrule
H6 chính & \path{wp9-confirmatory-result-2026-08-23.md}; biên bản từng ô B1/C1 &
\path{recheck_h6_frontier.py}; \path{h6-raw-movement-recheck.json}\\[4pt]
Thang trần & Biên bản từng ô trong \path{frontier-v1}; kết quả hai cửa sổ &
\path{recheck_wide_service.py}; \path{service-budget-per-cell.csv}\\[4pt]
Tầng 2 & \path{PREREGISTRATION.md}; \path{predictions.csv}; kết quả phân loại cuối &
\path{classify.py}; \path{tier2-recomputed-report.txt}\\[4pt]
W07, W07X & \path{W07-REPORT.md}; \path{W07X-PLAN.md}; \path{W07X-REPORT.md} &
\path{w07x_analyze.py}; \path{w07x_mechanism.py}\\[4pt]
T38 & \path{T38-PLAN.md}; \path{split-t38-v1.json}; \path{GAMMA-t38-v1.json} &
\path{recheck_t38.py}; \path{check_trip_overlap.py}\\
\bottomrule
\end{tabularx}
\end{table}

Một lần chạy cần được ghim vào đầu vào chuẩn hóa, cấu hình chính sách, phiên
bản runner và phiên bản adapter. Mã nguồn nhánh nghiên cứu được đọc tại
commit \texttt{24aa747}; đây là mốc của hồ sơ nghiên cứu mới nhất dùng trong
khóa luận, không thay thế các bản ghim riêng của H6 và tầng 2. Các kết quả
cũ vẫn gắn với runner cũ. Không xây lại runner đã đóng băng rồi dùng cùng
tên để đọc kết quả của một phiên bản khác.

Với H6, dữ liệu gốc nằm trong các cây kết quả
\path{wp9/confirmatory-h6-panela-v2} và \path{wp9/frontier-v1}. Mỗi ô giữ
\path{transcript-00.ndjson} và \path{run-00.json}. Cần đọc lịch sử lời hứa
ở điểm kiểm tra cuối, giữ nguyên các số mili giây, rồi cộng đúng tọa độ
còn mở. Tỉ số trung tâm là tỉ số \emph{hai tổng trên toàn panel}; lấy trung
bình các tỉ số từng khách sẽ trả lời một câu hỏi khác.

Với T38, cây \path{research/tier3-t38-v1} giữ cả các lần lỗi hạ tầng. Phép
tính dùng đúng 15 lần hiệu chỉnh và 170 lần kiểm tra hoàn tất đã chỉ định,
không cộng lần thử lại vào cỡ mẫu. Kiểm giao nhau phải nối
\path{selection-frame.json} qua \path{normalizer-configuration.json} tới
CSV nguồn, rồi dùng định danh chuyến gốc. Băm của tệp chia, các CSV nguồn
và danh sách chuyến giao nhau được giữ trong
\path{t38-actual-trip-id-overlap.json}.

Các hình được sinh bằng cùng các CSV/JSON đã kiểm, xuất PDF và SVG vector.
Tệp \path{build_verified_figures.py} dựng bốn hình kết quả;
\path{mechanism_figure.py} dựng ví dụ ba mốc và không được dùng như dữ
liệu thực nghiệm. Gói biên tập giữ mã phân tích, bảng từng ô và hồ sơ kiểm
trong thư mục \path{audit}. Tính lại một con số từ cùng biên bản kiểm phép
trích xuất; nó không tạo thêm một mẫu thực nghiệm độc lập.
